\documentclass{beamer}
% Computer Science 1 — shared Beamer preamble.
% \input this file after \documentclass{beamer} in every Monday lecture
% deck so all decks share one visual system. Do not fork per-week copies of
% this file — edit it once here and every deck picks up the change on next
% compile. Vendored from professional_minds/presentations/beamer/theme/
% preamble.tex (same pattern semester_kickoff_week's own vendored copy
% uses), kept local rather than \input across repos so this repo stays
% self-contained and usable without a professional_minds checkout.

\usetheme{Boadilla}
\usecolortheme{seahorse}
\setbeamertemplate{navigation symbols}{}
\setbeamertemplate{footline}[frame number]

% Speaker notes: \note{...} on any frame. Compiling with the default
% `pdflatex deck.tex` produces a slides-only PDF (notes are embedded but
% hidden). Compiling with
% `pdflatex "\PassOptionsToPackage{show}{pgfpages}\input{deck.tex}"`
% or uncommenting the line below produces an instructor copy with each
% slide's notes on a second page — see README.md for both commands.
% \setbeameroption{show notes on second screen=right}

\usepackage{pgfpages}

\newcommand{\SessionTitleSlide}[4]{%
  % #1 = Week/Day (e.g. "Week 2 -- Monday"), #2 = session title,
  % #3 = essential question, #4 = primary source citation
  \begin{frame}
    \titlepage
  \end{frame}
  \begin{frame}{Essential Question}
    \Large #3

    \vspace{1em}
    \normalsize\emph{#4}
  \end{frame}
}

\usepackage{hyperref}

\title{Week 11 Monday}
\subtitle{Dictionaries: A Dictionary of Dictionaries}
\author{CS1}
\date{}

\begin{document}

\SessionTitleSlide{Week 11 -- Monday}{Dictionaries: Look It Up by Key}{How do we find one thing by name, and can a dictionary hold other dictionaries?}{Week 11 plan; \texttt{lessons/07-collections.md}; AI Fluency Lens 11 (Verify) tie-in}

\begin{frame}[fragile]{Rock-Paper-Scissors: Rules as Data}
\begin{verbatim}
rules = {
 "rock": {"rock": "tie", "paper": "lose",
          "scissors": "win"},
 "paper": {"rock": "win", "paper": "tie",
           "scissors": "lose"},
 "scissors": {"rock": "lose", "paper": "win",
              "scissors": "tie"},
}
print(rules["rock"]["scissors"])  # win
\end{verbatim}
\begin{itemize}
    \item A dictionary whose values are dictionaries: \texttt{rules[player][opponent]}.
    \item First key picks the row, second key picks the column. Two lookups, one answer.
    \item Predict first: what is \texttt{rules["paper"]["rock"]}?
\end{itemize}
\note{Ask for predictions before running. Then ask: how many if statements would the same game need? Nine. Changing a rule here means changing data, not logic. Run the file live: python3 rps\_dict\_of\_dicts.py}
\end{frame}

\begin{frame}[fragile]{Walk It, Score It, Guard It}
\begin{verbatim}
for mine in rules:            # outer keys
    for theirs in rules[mine]:  # inner keys
        print(mine, theirs, rules[mine][theirs])

score = {"win": 0, "lose": 0, "tie": 0}
score[judge("rock", "scissors")] += 1
\end{verbatim}
\begin{itemize}
    \item A second dictionary keeps score: keys are results, values are counts.
    \item A move that is not a key raises \texttt{KeyError}; check with \texttt{in} first.
    \item \textbf{Verify (Lens 11):} try an empty string, a capital letter, a move not in the table.
    \item \textbf{Run it, read the walkthrough:} \href{https://swosu.instructure.com/courses/74029/pages/week-11-overview-lists-and-dictionaries-continued}{Canvas, Week 11 Overview page, ``Rock-Paper-Scissors example''}
\end{itemize}
\note{Code and walkthrough are linked from that Canvas page (Python file plus README). Decide/Compare tie-in: why a dict of dicts and not nine ifs, a list of lists, or one dict keyed by a pair? What changes when we add lizard and spock?}
\end{frame}

\end{document}
